\chapter{Revue de la littérature}
\label{chap:revue}

Ce chapitre situe la présente étude dans son double contexte intellectuel : d'une part, les
travaux relatifs à la segmentation clientèle par le modèle RFM et à son usage dans les
plateformes de gestion ; d'autre part, les méthodes d'évaluation par simulation permettant
de mesurer la conformité et la robustesse d'un module algorithmique. La revue s'organise en
quatre sections. La première ancre la problématique dans la réalité des petits commerçants
congolais et dans les limites des solutions SaaS disponibles. La deuxième présente le modèle
RFM en distinguant rigoureusement ses trois familles d'implémentation, leurs hypothèses
respectives et leurs limites documentées. La troisième pose les critères généraux permettant
d'évaluer un module de segmentation, sans anticiper les résultats de la campagne. La quatrième
introduit le cadre méthodologique de la simulation par Monte-Carlo, en particulier le cadre
ADEMP proposé par \citet{Morris2019}, et identifie les configurations dégénérées susceptibles
d'affecter la stabilité des segmentations.

% ------------------------------------------------------------------
\section{Artefact et besoin : SaaS de gestion et segmentation des petits commerçants}
\label{sec:revue-artefact}

\subsection{Le petit commerce congolais et le déficit d'outils adaptés}
\label{subsec:revue-contexte-rdc}

La République démocratique du Congo compte une part importante de son activité économique
dans le secteur des très petites entreprises et du commerce informel. Ces acteurs exercent
souvent sans comptabilité formelle, dans un environnement monétaire dual (franc congolais et
dollar américain), avec un accès irrégulier à l'infrastructure numérique. \citet{Yoseph2023}
documentent la difficulté structurelle des TPE africaines à adopter des systèmes de gestion
conçus pour des contextes économiques distincts : les solutions génériques, développées pour
des marchés occidentaux ou asiatiques, intègrent rarement les logiques de devises multiples,
de transactions partielles ou de clientèle non bancarisée.

Dans ce contexte, ELMES-TEKA est conçu comme un SaaS de gestion commerciale multi-tenant,
destiné à fournir aux entrepreneurs congolais formels et informels un outil de suivi des
ventes, des stocks, des dépenses et des clients. L'une de ses fonctionnalités analytiques
est un module de segmentation clientèle fondé sur le modèle RFM, dont l'objectif est de
classer les acheteurs selon leur valeur comportementale afin d'orienter les décisions
commerciales.

\subsection{La segmentation comme fonction analytique dans un SaaS de gestion}
\label{subsec:revue-segmentation-saas}

\citet{Salminen2023} soulignent que les modules analytiques intégrés dans des plateformes
commerciales répondent à une logique différente de celle des outils analytiques autonomes :
ils doivent produire des résultats stables, comparables d'un exercice à l'autre, et
interprétables par des utilisateurs non statisticiens. Dans cette optique, un module de
segmentation doit non seulement produire des résultats corrects, mais également se comporter
de manière prévisible face à des variations mineures des données d'entrée.

La question de la stabilité est particulièrement sensible dans le contexte des petits
commerçants, dont les données transactionnelles sont souvent bruitées : erreurs de saisie,
arrondis monétaires, doublons de commandes, corrections a posteriori. Un module incapable de
traiter ces imperfections de manière cohérente produirait des classements instables, minant
la confiance des utilisateurs dans les recommandations commerciales générées.

\subsection{Positionnement de la présente étude}
\label{subsec:revue-positionnement}

L'évaluation conduite dans ce mémoire relève de la démarche de Design Science Research
(\emph{DSR}). Dans ce cadre, la fonction \texttt{computeRFM} est l'artefact dont l'étape
\emph{evaluate} consiste à mesurer la conformité à une référence indépendante et la
robustesse aux perturbations réalistes. Cette posture épistémologique, dite pragmatique et
post-positiviste, est cohérente avec l'objectif d'amélioration d'un produit en cours de
développement : il ne s'agit pas de tester une hypothèse théorique sur la segmentation en
général, mais d'établir des garanties mesurables sur un module logiciel précis, dans les
conditions d'usage définies par son contexte.

% ------------------------------------------------------------------
\section{Le modèle RFM : données, quantiles, familles et limites connues}
\label{sec:revue-rfm}

\subsection{Origines et composantes du modèle}
\label{subsec:revue-rfm-origines}

Le modèle RFM est fondé sur trois dimensions comportementales d'un client :
la \emph{Récence} ($R$), soit le temps écoulé depuis le dernier achat ; la
\emph{Fréquence} ($F$), soit le nombre de transactions sur une période donnée ; et le
\emph{Montant} ($M$, parfois désigné \emph{Monetary Value}), soit la somme totale dépensée
sur la même période. Ces trois dimensions sont calculées individuellement pour chaque client
à partir de l'historique transactionnel disponible.

\citet{Fader2005} attribuent l'origine pratique de ce modèle aux praticiens du marketing
direct des années 1990, qui cherchaient un outil simple pour hiérarchiser les listes de
clients en vue de campagnes postales. La force du modèle réside précisément dans cette
simplicité : trois valeurs numériques suffisent à résumer le profil comportemental d'un
client, sans nécessiter de modélisation probabiliste complexe ni de données démographiques.

Dans ELMES-TEKA, la Récence est définie comme le nombre de jours depuis le dernier achat
enregistré pour un exercice donné ; la Fréquence comme le nombre de commandes acceptées
(statuts \texttt{CONFIRMED}, \texttt{PAID} ou \texttt{DELIVERED}) ; le Montant comme la
somme des montants de ces mêmes commandes.

\subsection{Les trois familles d'implémentation RFM}
\label{subsec:revue-rfm-familles}

Une confusion fréquente dans la littérature consiste à traiter le modèle RFM comme une
méthode uniforme. Il existe en réalité trois familles d'implémentation distinctes, reposant
sur des hypothèses et des propriétés mathématiques différentes.

\paragraph{Famille 1 — Score par règles et terciles (ou quintiles).}
La première famille, et la plus répandue dans les outils commerciaux, attribue à chaque
client un score discret par découpage en quantiles de la distribution empirique de chaque
dimension. \citet{Salminen2023} décrivent ce mécanisme comme une normalisation ordinale :
pour chaque dimension, les seuils sont calculés de manière à répartir les clients en groupes
d'effectifs approximativement égaux. Un score de tercile $s \in \{1, 2, 3\}$ est attribué en
fonction du rang dans la distribution. C'est cette famille qu'implémente ELMES-TEKA, avec
cinq catégories finales (\emph{Champions}, \emph{Loyaux}, \emph{À risque},
\emph{Petits acheteurs}, \emph{Potentiels}) définies par des règles portant sur les scores
$s_R$, $s_F$ et $s_M$.

La simplicité de cette approche a un coût bien documenté. \citet{ViannaFilho2026} rappellent
que le découpage en terciles est sensible aux égalités (\emph{ties}) : lorsque plusieurs
clients présentent exactement la même valeur sur une dimension, leur affectation à un groupe
dépend du mécanisme de résolution des ex-æquo, qui peut varier d'une implémentation à
l'autre. Cette instabilité est particulièrement marquée pour les corpus de petite taille, où
un seul client peut représenter un tiers de la distribution. Par ailleurs, \citet{Sobantu2025}
soulignent que la méthode par terciles n'est pas monotone au sens strict pour les valeurs
situées précisément aux seuils : une variation infinitésimale d'un montant peut entraîner un
changement de score si elle déplace le client d'un côté ou de l'autre d'un seuil.

\paragraph{Famille 2 — Clustering non supervisé (K-means et variantes).}
La deuxième famille utilise des algorithmes de regroupement pour partitionner l'espace
tridimensionnel $(R, F, M)$ en clusters homogènes. \citet{Rajendran2024} proposent une
implémentation fondée sur K-means pour la segmentation de clientèle dans le commerce de
détail, en justifiant le nombre de clusters par le critère de la silhouette ou de la
méthode du coude. \citet{Sobantu2025} comparent K-means et K-médoïdes dans un contexte de
commerce en ligne africain et concluent que K-médoïdes est plus robuste aux valeurs
aberrantes, au prix d'une complexité computationnelle plus élevée.

Cette famille se distingue de la première par son caractère non déterministe au sens de
la graine d'initialisation et par l'absence de catégories sémantiques prédéfinies : les
segments obtenus dépendent des données et ne correspondent pas nécessairement aux cinq
groupes habituels du modèle RFM classique.

\paragraph{Famille 3 — Modèles stochastiques (Pareto/NBD et dérivés).}
La troisième famille modélise explicitement le processus d'achat comme un processus
stochastique. \citet{Fader2005} proposent le modèle Pareto/NBD (\emph{Negative Binomial
Distribution}), dans lequel la fréquence des achats suit une loi de Poisson dont le paramètre
varie entre clients selon une loi gamma, et la durée de vie active du client suit une loi
exponentielle dont le paramètre varie selon une loi gamma. Ce modèle permet d'estimer la
probabilité qu'un client soit encore «~vivant~» et de prédire son comportement futur.

La complexité de calibration et d'interprétation de cette famille la rend peu adaptée aux
outils grand public ou aux contextes de petits commerces disposant de données limitées.
ELMES-TEKA n'appartient pas à cette famille.

\subsection{Limites communes et problème des catégories dégénérées}
\label{subsec:revue-rfm-limites}

Quelle que soit la famille d'implémentation, le modèle RFM présente des limites structurelles
reconnues. \citet{Salminen2023} rappellent que le modèle ne capture pas les motivations
d'achat, la fidélité attitudinale, ni les effets de saisonnalité à court terme. La dimension
Montant ignore la composition du panier et peut conduire à ranger dans la même catégorie
des clients aux profils très différents.

Un cas limite particulièrement important pour l'évaluation d'un module de segmentation est
celui des corpus dits \emph{dégénérés} ou \emph{ex-æquo} : un corpus où tous les clients
présentent les mêmes valeurs sur une ou plusieurs dimensions. Dans ce cas, le découpage en
terciles est mathématiquement sous-déterminé, car toute partition en groupes d'effectifs
égaux est valide. \citet{ViannaFilho2026} identifient ce cas comme une configuration
hors-domaine pour les méthodes par quantiles, dans la mesure où les hypothèses implicites
de la méthode — la dispersion suffisante des valeurs pour définir des seuils stables —
ne sont pas satisfaites. Les comportements observés dans ce type de corpus ne peuvent pas
être interprétés comme des mesures de performance au sens habituel ; ils décrivent un
comportement hors-domaine défini par la construction même du corpus.

% ------------------------------------------------------------------
\section{Évaluer un module de segmentation : conformité, robustesse et reproductibilité}
\label{sec:revue-evaluation}

\subsection{Conformité à une référence indépendante}
\label{subsec:revue-conformite}

L'évaluation d'un module algorithmique peut poursuivre deux objectifs distincts. Le premier
est la vérification de la conformité (\emph{verification}) : le module produit-il les sorties
attendues pour des entrées données ? Cette question suppose l'existence d'une référence
indépendante, par rapport à laquelle les sorties du module peuvent être comparées. Dans le
cas d'un module de segmentation basé sur des terciles, la référence peut être une
implémentation dans un langage ou un environnement différent, en s'assurant que les
hypothèses sur la représentation des nombres sont explicitement contrôlées.

\citet{Salminen2023} insistent sur l'importance de définir précisément ce que l'on entend
par «~sortie correcte~» avant de mesurer la conformité : pour un module de segmentation
numérique, les erreurs d'arrondi et les accumulations de virgule flottante peuvent produire
des sorties différant de la référence sans qu'aucune erreur logique ne soit en cause.
L'unité d'analyse doit donc être clairement posée — l'essai, le client, la dimension — et
le critère de non-conformité doit être explicitement défini avant la collecte des données.

\subsection{Robustesse aux perturbations}
\label{subsec:revue-robustesse}

Le second objectif est la vérification de la robustesse (\emph{validation}) : le module
produit-il des résultats stables lorsque les données d'entrée sont soumises à des
perturbations réalistes ? Cette question est distincte de la précédente : un module peut
être conforme à sa référence et néanmoins instable, ou inversement.

Dans le contexte d'une plateforme de gestion commerciale, les perturbations réalistes
incluent les erreurs de saisie monétaire (arrondis, corrections), les modifications de dates
de transaction, les doublons de commandes et les suppressions accidentelles. La robustesse
d'un module de segmentation peut être mesurée par le taux de changement de segment parmi
les clients communs à deux partitions produites à partir d'un corpus initial et d'un corpus
perturbé. \citet{Rajendran2024} utilisent une logique analogue dans l'évaluation de leur
méthode de segmentation, en mesurant la stabilité des clusters sous rééchantillonnage.

La définition du critère de compatibilité requiert de poser une borne acceptée en amont.
\citet{Lakens2013} soulignent l'importance de distinguer la significativité statistique de
la pertinence pratique : un test qui détecte un changement statistiquement significatif peut
décrire un effet trop faible pour avoir une conséquence commerciale. Dans une étude de
robustesse, la borne opérationnelle — le seuil au-delà duquel un taux de changement est
considéré comme incompatible — doit être définie par des considérations métier, non par un
seuil statistique arbitraire. La borne supérieure d'un intervalle de confiance Monte-Carlo
constitue un instrument de décision naturel pour ce type de comparaison.

\subsection{Perspectives analytiques complémentaires}
\label{subsec:revue-perspectives}

Plusieurs outils analytiques sont mobilisables dans une phase ultérieure d'évaluation plus
approfondie, lorsqu'une cible indépendante est disponible. L'analyse de variance (ANOVA),
le test de Welch adapté aux variances inégales, les tests post-hoc de comparaisons multiples,
les mesures d'effet (\emph{d} de Cohen, $\eta^2$), la régression logistique pour la
prédiction du changement de segment, et les indicateurs de discrimination tels que la courbe
ROC et l'AUC permettent de caractériser les effets de manière plus fine. Ces méthodes
supposent cependant la disponibilité d'une variable cible indépendante du processus de
segmentation lui-même, ce qui n'est pas le cas dans une étude de robustesse interne.
Elles constituent des perspectives naturelles pour une évaluation comparative entre plusieurs
implémentations ou sur des données de terrain réelles, mais elles ne font pas partie du
protocole de la campagne présentée dans ce mémoire.

La reproductibilité des résultats est une condition supplémentaire, distincte de la
conformité et de la robustesse. Elle concerne la capacité d'un tiers à reconstruire
identiquement les résultats à partir du code source, des données et de la procédure
documentés. Dans le cadre de cette étude, la reproductibilité est assurée par la fixation
de graines aléatoires, la documentation des versions logicielles et l'archivage complet de
l'environnement d'exécution.

% ------------------------------------------------------------------
\section{Simulation pour l'évaluation de méthodes : ADEMP, Monte-Carlo et catégories dégénérées}
\label{sec:revue-simulation}

\subsection{Le rôle de la simulation dans l'évaluation des méthodes statistiques}
\label{subsec:revue-simulation-role}

L'évaluation d'une méthode par simulation est une pratique établie en statistique et en
informatique lorsque les propriétés théoriques de la méthode sont difficiles à dériver
analytiquement, ou lorsque les hypothèses requises par une démonstration formelle ne sont
pas satisfaites dans le cas d'usage. \citet{Ribatet_nd} présentent la simulation de
Monte-Carlo comme un outil générique permettant d'approximer des quantités mathématiques
par la répétition d'un grand nombre d'expériences aléatoires indépendantes, dont la
moyenne converge vers la quantité d'intérêt par la loi des grands nombres.

Dans le contexte de l'évaluation d'un module logiciel, la simulation sert un double objectif.
D'une part, elle permet d'explorer le comportement du module sur un espace de configurations
plus large que ce que permettrait l'observation directe. D'autre part, elle fournit des
estimations assorties d'intervalles de confiance, permettant de quantifier l'incertitude
liée au nombre fini de répétitions.

\subsection{Le cadre ADEMP}
\label{subsec:revue-ademp}

\citet{Morris2019} proposent un cadre de présentation des études de simulation, désigné
par l'acronyme ADEMP, structurant cinq composantes essentielles d'une étude de simulation :

\begin{enumerate}
  \item \textbf{Aims} (\emph{objectifs}) : l'objectif scientifique de l'étude, posé sous
    forme de question ou d'hypothèse mesurable ;
  \item \textbf{Data-generating mechanisms} (\emph{mécanismes générateurs de données}) :
    les procédures permettant de construire les jeux de données synthétiques soumis à la
    méthode évaluée ;
  \item \textbf{Estimands} (\emph{cibles}) : les quantités que l'étude cherche à estimer ou
    à évaluer, définies indépendamment de la méthode ;
  \item \textbf{Methods} (\emph{méthodes}) : la ou les méthodes comparées ou évaluées, dont
    les propriétés sont mesurées par rapport aux cibles ;
  \item \textbf{Performance measures} (\emph{mesures de performance}) : les indicateurs
    quantitatifs permettant de répondre à l'objectif, accompagnés de leur incertitude Monte-Carlo.
\end{enumerate}

Ce cadre est mobilisé dans la présente étude pour présenter la campagne de simulation
conduite sur la fonction \texttt{computeRFM}. Son utilisation comme grille de lecture
n'implique pas un préenregistrement historique de la campagne ; il s'agit d'un outil
de structuration de la présentation, adopté \emph{a posteriori} pour garantir la clarté
et la reproductibilité du rapport d'évaluation.

\citet{Morris2019} insistent par ailleurs sur la nécessité de quantifier l'incertitude
Monte-Carlo des mesures de performance elles-mêmes. Pour un estimateur de proportion tel
que le taux de non-conformité ou le taux de changement de segment, l'incertitude est
estimée par un intervalle de confiance fondé sur la loi $t$ à partir de la variance
empirique des répétitions. Cette précision méthodologique est indispensable pour éviter de
conclure à la robustesse d'un module sur la seule base d'un taux moyen, sans quantifier
la variabilité inter-répétitions.

\subsection{Mécanismes générateurs et profils de corpus}
\label{subsec:revue-simulation-dgm}

La conception des mécanismes générateurs de données est une étape critique d'une étude de
simulation. \citet{Morris2019} distinguent deux types de mécanismes : les mécanismes
\emph{paramétriques}, qui construisent des données selon un modèle probabiliste explicite,
et les mécanismes \emph{empiriques}, qui rééchantillonnent des données observées. Dans le
cas présent, des mécanismes paramétriques sont utilisés pour produire des corpus synthétiques
de clients caractérisés par leur Récence, leur Fréquence et leur Montant.

Trois profils de distribution sont pertinents pour cette étude :

\begin{itemize}
  \item Le profil \emph{régulier} (\texttt{regular}), dans lequel les valeurs sont
    distribuées de manière relativement homogène, reproduisant un contexte commercial
    ordinaire avec une certaine diversité de comportements clients ;
  \item Le profil \emph{dissymétrique} (\texttt{skewed}), dans lequel quelques clients
    concentrent une part importante de l'activité, reproduisant une distribution de Pareto
    courante dans le commerce de détail ;
  \item Le profil \emph{ex-æquo} (\texttt{ties}), dans lequel tous les clients présentent
    exactement les mêmes valeurs sur les trois dimensions, constituant un cas dégénéré au
    sens de la méthode des terciles.
\end{itemize}

Ce dernier profil mérite une attention particulière. Comme le montrent \citet{ViannaFilho2026},
une segmentation par terciles sur un corpus où tous les clients sont identiques est
mathématiquement ambiguë : les seuils de terciles ne peuvent pas être définis de manière
unique, et la résolution des ex-æquo par tout algorithme concret introduit une forme
d'arbitraire. Les résultats obtenus sur ce profil ne mesurent pas la robustesse au sens
de la résistance aux perturbations, mais décrivent le comportement de l'algorithme dans une
configuration que son domaine de validité n'inclut pas. Ils constituent néanmoins une
information utile pour circonscrire les limites du module évalué.

\subsection{Perturbations et cellules de la campagne}
\label{subsec:revue-simulation-perturbations}

La conception des perturbations appliquées aux corpus reflète les conditions d'usage
réelles de la plateforme. \citet{Morris2019} recommandent que les perturbations soient
motivées par le contexte applicatif plutôt que par des choix purement statistiques : elles
doivent correspondre à des erreurs ou des variations plausibles dans les données de terrain.

Dans le contexte d'une plateforme de gestion commerciale pour petits commerçants congolais,
les perturbations réalistes incluent :

\begin{itemize}
  \item le bruit monétaire, c'est-à-dire les variations de montant liées aux arrondis de
    prix ou aux corrections de saisie ;
  \item les décalages de dates, introduits par des corrections de dates de transaction ou
    des enregistrements différés ;
  \item les doublons de commandes, résultant d'enregistrements redondants ;
  \item les commandes manquantes, résultant de suppressions ou de rejets a posteriori.
\end{itemize}

Pour chaque type de perturbation, plusieurs intensités sont testées afin de caractériser la
relation entre l'intensité de la perturbation et le taux de changement de segment. La
conception en cellules — croisement de la taille du corpus, du profil de distribution et
du scénario de perturbation — permet d'isoler les effets de chaque facteur et d'identifier
les configurations pour lesquelles le module est le plus sensible.

\subsection{Intervalles de confiance Monte-Carlo et critère de décision}
\label{subsec:revue-simulation-ic}

L'incertitude des estimateurs de proportion calculés par simulation est quantifiable par
des intervalles de confiance Monte-Carlo. Pour un taux observé $\hat{p}$ sur $B$ répétitions,
l'intervalle de confiance asymptotique à $(1-\alpha)$ fondé sur la loi normale est :

\begin{equation}
  \hat{p} \;\pm\; z_{1-\alpha/2} \;\sqrt{\frac{\hat{p}(1-\hat{p})}{B}}
  \label{eq:ic-monte-carlo}
\end{equation}

\noindent où $z_{1-\alpha/2}$ est le quantile de la loi normale standard. Pour un nombre
modéré de répétitions, \citet{Morris2019} recommandent d'utiliser la loi $t$ à $B - 1$ degrés
de liberté à la place de la loi normale, ce qui produit des intervalles légèrement plus
conservateurs mais mieux calibrés.

La décision de compatibilité repose sur la comparaison entre la borne supérieure de
l'intervalle de confiance et une borne opérationnelle définie avant la campagne. Ce
mécanisme est cohérent avec la recommandation de \citet{Lakens2013} de recourir à des
tests d'équivalence ou d'infériorité pour les études de robustesse, dans lesquelles
l'objectif est de montrer qu'un effet est suffisamment petit pour être négligeable, et non
qu'il est différent de zéro.

% ------------------------------------------------------------------
\section{Synthèse et lacunes identifiées}
\label{sec:revue-synthese}

La revue de la littérature fait apparaître les éléments suivants, qui structurent la
contribution de la présente étude.

Premièrement, le modèle RFM par terciles, tel qu'implémenté dans ELMES-TEKA, est une
approche légère et interprétable, mais qui comporte des limites structurelles bien
documentées : sensibilité aux ex-æquo, instabilité aux seuils, indifférence aux profils
monetaires complexes. Ces limites sont inhérentes à la famille d'implémentation choisie
et ne constituent pas nécessairement un défaut de l'artefact, mais elles délimitent son
domaine de validité.

Deuxièmement, l'évaluation d'un module de segmentation dans un SaaS de gestion doit
distinguer la conformité à une référence — question de vérification — et la robustesse
aux perturbations — question de validation. Ces deux dimensions requièrent des protocoles
distincts, des critères de décision définis avant la collecte et une quantification explicite
de l'incertitude.

Troisièmement, la simulation de Monte-Carlo, encadrée par le schéma ADEMP de
\citet{Morris2019}, fournit un protocole méthodologique adapté à l'évaluation d'un module
logiciel lorsque les propriétés analytiques de la méthode sont difficiles à dériver. Elle
permet de couvrir un espace large de configurations, d'isoler les effets des facteurs
d'intérêt et de produire des estimateurs accompagnés de leur incertitude.

Enfin, la littérature consultée ne documente pas d'étude d'évaluation par simulation du
module RFM d'un SaaS de gestion commerciale destiné aux petits commerçants d'Afrique
subsaharienne. La présente étude comble partiellement cette lacune, en produisant des
mesures de conformité et de robustesse sur un artefact réel, dans un contexte applicatif
précis.
